\subsection{The exact first seed cannot be repaired by an $S$-unit}

\begin{statement}{A.159}{Proved}
{A finite unit certificate excludes every $S$-unit continuation of the specified seed}
Retain the actual fields, embeddings and elements of A.156, and put
\[
 S=\{2,3,11,59,\infty\},\qquad
 E_i=\mathcal O_{K_i,S_i}^{\times},\qquad
 \widehat E_i^{(3)}=\varprojlim_r E_i/E_i^{3^r}.
\tag{A.159.1}
\]
Let \(K_3=\Q(\rho_3)\), where \(\rho_3^3=\rho_2=d^2t\).
There is no element \(b_2\in\widehat E_2^{(3)}\) satisfying both
\[
 N_{K_2/K_1}b_2=1-\alpha,
 \qquad
 b_2\in N_{K_3/K_2}\bigl(\widehat{K_3^{\times}}^{(3)}\bigr).
\tag{A.159.2}
\]
In particular no actual $S$-unit preimage of the exact seed
\(\beta_1=1-\alpha\) extends to the next field norm.
The assertion allows every $S$-unit modification of A.156's \(\beta_2\).
It leaves open norm families whose first coordinate is a different element
of the norm fibre over \(-2\), and modifications with additional support.
Moreover \(\operatorname{Cl}(K_2(\zeta_3))\{3\}\) contains an
explicit subgroup isomorphic to \((\Z/3)^2\), generated by the
three-primary differences of the two conjugate quadratic primes over
eleven and fifty-nine described in the companion proof.
\end{statement}

\begin{proof}
Write \(U=\alpha^2-2\); its inverse is
\(4+3\alpha+2\alpha^2\), and its rational norm is one.
Let \(V^+\) be the group of global units of \(K_2\) of rational norm one.
The field has signature \((1,4)\), so Dirichlet's theorem makes \(V^+\)
a free abelian group of rank four.  Its torsion is trivial: \(K_2\) is
contained in \(\R\), and the odd degree makes the norm of \(-1\) equal
\(-1\).  More explicitly,
\(\mathcal O_{K_2}^{\times}=\{\pm1\}\times V^+\),
so its pro-three completion is exactly \(V^+\otimes\Z_3\).

Here is the finite exact certificate used in the proof.  Four units
\(u_1,u_2,u_3,u_4\in V^+\), with \(u_1=U\), satisfy
\[
 N_{K_2/K_1}u_i=U^{k_i},\qquad
 (k_1,k_2,k_3,k_4)=(3,57,-3,1614).
\tag{A.159.3}
\]
The complete coefficient supplement for this statement gives literal
integer coordinates for all four units in the elements
\(e_j=t^j/d^{\lfloor2j/3\rfloor}\).  Each \(e_j\) is integral,
since \(e_j^9=3^j d^{6j-9\lfloor2j/3\rfloor}\) is an integer.
For \(u=A+Bt+Ct^2\), with \(A,B,C\in K_1\) and \(w=d^2\alpha\),
the exact relative-norm formula is
\[
                 Nu=A^3+wB^3+w^2C^3-3wABC.
\]
Substitution of the displayed integer coordinates gives A.159.3.
Since \(N_{K_1/\Q}U=1\), all four integral elements have absolute
norm one and are units.  The complete rational coordinates and an
additional monic-characteristic-polynomial check are retained in
\path{code/two_inert_relative_unit_repair_certificate_data.json}.
The four rows of the following table are exact residue evaluations at
unramified degree-one primes of \(K_2\).  The entry \(r\) is the residue
of \(t\), and \(\omega\) is the chosen nonidentity cubic root of unity
in \(\F_p\):
\[
\begin{array}{c|r|r|rrrr|rrrr}
 p&r&\omega&\overline u_1&\overline u_2&\overline u_3&\overline u_4
 &\lambda_1&\lambda_2&\lambda_3&\lambda_4\\\hline
61&27&47&23&23&4&42&0&0&2&2\\
103&40&56&62&82&62&86&2&1&2&1\\
151&33&32&114&96&105&120&1&2&2&1\\
193&8&84&179&167&70&16&0&1&1&1
\end{array}
\tag{A.159.4}
\]
Here \(r^9=3d^6\pmod p\) and
\(\overline u_i^{(p-1)/3}=\omega^{\lambda_i}\).
The matrix of the last four columns has integer determinant two.
All denominators in the unit coordinates are invertible at these primes,
so the evaluations are literal ring homomorphisms.  Consequently a product
\(\prod_i u_i^{e_i}\) can be a cube in \(K_2\) only if every
\(e_i\) is divisible by three.

For clarity, this proves the required saturation without an analytic
class-number or regulator calculation.  It first proves that the four
units are multiplicatively independent: divide any nonzero integer
relation by its highest common power of three and use the determinant.
Thus \(H=\langle u_1,u_2,u_3,u_4\rangle\) has finite index in \(V^+\).
If that index were divisible by three, an element \(v\notin H\) with
\(v^3\in H\) would give such a cube relation.  The determinant would
write \(v^3=h^3\) with \(h\in H\); absence of \(\mu_3\) in
\(K_2\) gives \(v=h\), a contradiction.  Therefore
\[
 [V^+:H]\ \text{is prime to three},\qquad
 H\otimes\Z_3\xrightarrow{\sim}V^+\otimes\Z_3.
\tag{A.159.5}
\]

Let \(Q_{11}\) and \(Q_{59}\) denote the quadratic-residue places of
\(K_2\) used in A.156.  Their residue fields have order \(q^2\).
Define the continuous cubic residue characters on local units by
\[
 \lambda_q(v)=\overline v^{(q^2-1)/3}\in\mu_3(\F_{q^2}).
\tag{A.159.6}
\]
Exact reduction of the same four candidates gives
\[
 v_i=u_iU^{-k_i/3},\qquad N_{K_2/K_1}v_i=1,
 \qquad \lambda_{11}(v_i)=\lambda_{59}(v_i)=1
 \quad(1\leq i\leq4).
\tag{A.159.7}
\]
This check uses the same explicit finite fields as A.156.12--14.
In a power-basis expression \(A_0(\alpha)+A_1(\alpha)t+
A_2(\alpha)t^2\), the certificate checks that \(A_0,A_1\) are integral
at both primes and that the denominators of \(A_2\) have valuation at
most one.  Since \(v(t)=2\) and \(v(q)=3\), the reduction is precisely
\(A_0(\overline\alpha)\).

Every completed global unit of relative norm one has a unique expression
\(v=\prod_i u_i^{z_i}\), with \(z_i\in\Z_3\), by A.159.5.
Its norm is \(U^{\sum_i k_i z_i}\).  The image of the infinite-order
unit \(U\) has infinite order in the pro-three completion of the
finitely generated unit group of \(K_1\), so
\(\sum_i k_i z_i=0\).  Division by three in the torsion-free ring
\(\Z_3\) then gives
\[
 v=\prod_i v_i^{z_i},\qquad
 \lambda_{11}(v)=\lambda_{59}(v)=1.
\tag{A.159.8}
\]
The same conclusion for actual units also follows by raising to the
prime-to-three index in A.159.5.

We now prove the precise reduction from $S$-units to global units.
At each old place over eleven or fifty-nine, the equation
\(t^3=d^2\alpha\) has radicand valuation two.  It therefore defines
one totally ramified cubic extension.  At three the radicand has
valuation one, again giving one totally ramified cubic extension.
The relative ideal norm at each of these places multiplies valuations
by the residue degree one.  Hence a relative norm-one $S$-unit has
valuation zero at every place over \(3,11,59\).

At two, the actual factorization and the map on residues are
\[
 T^9-1=(T-1)(T^2+T+1)(T^6+T^3+1),\qquad
 \overline\alpha=\overline t^{3}.
\tag{A.159.9}
\]
All three factors are irreducible over \(\F_2\), of degrees one, two
and six.  Denote their primes by \(P_0,P_2,P_6\).
The first two lie over the old degree-one prime \(p_0\); the third
lies over the old degree-two prime \(p_2\).  The actual ideal-norm map is
\[
 P_0\longmapsto p_0,\qquad
 P_2\longmapsto2p_0,\qquad
 P_6\longmapsto3p_2.
\tag{A.159.10}
\]
It follows that the norm-zero divisor subgroup supported over two is
\(\Z(P_2-2P_0)\), and after completion it is
\(\Z_3(P_2-2P_0)\).
On the other hand A.156 gives
\[
 \operatorname{div}(\beta_1)=P_0+P_2,\quad
 \operatorname{div}(\beta_2)=P_0,\quad
 h=\frac{\beta_1}{\beta_2^3},\quad
 \operatorname{div}(h)=P_2-2P_0,\quad Nh=1.
\tag{A.159.11}
\]
The residue identity \(\overline\beta_2^{3}=
\overline\beta_1\) at both quadratic places gives
\(\lambda_{11}(h)=\lambda_{59}(h)=1\).

The divisor sequence for the finitely generated $S$-unit group has a
free abelian image.  It therefore splits as a sequence of abelian groups;
its pro-three completion has kernel exactly the completed global unit
group.  If \(z\in\widehat E_2^{(3)}\) has relative norm one, write
\(\operatorname{div}(z)=b(P_2-2P_0)\), \(b\in\Z_3\).
Then \(zh^{-b}\) is a completed global unit of relative norm one.
Equations A.159.8 and A.159.11 prove
\[
 \lambda_{11}(z)=\lambda_{59}(z)=1
 \quad\text{for every }z\in
 \ker\bigl(N:\widehat E_2^{(3)}\to\widehat E_1^{(3)}\bigr).
\tag{A.159.12}
\]
The inclusions of completed $S$-unit groups into completed field
multiplicative groups are injective: an $S$-unit that is a \(3^r\)-th
power in the field has an $S$-unit root, as its valuations outside $S$
are zero.

Finally suppose \(Nb_2=\beta_1\).  Then
\(b_2/\beta_2\) belongs to the kernel in A.159.12.  Thus \(b_2\)
has the same two nonidentity cubic characters as the explicit
\(\beta_2\) in A.156.  The next extension is a single totally ramified
cubic at each quadratic place, since \(v(\rho_2)=8\).
The norm of a local unit has cube residue; this remains true after
pro-three completion because the residue-character quotient is finite.
Any completed norm with valuation zero comes from a completed unit,
since the residue degree is one.  Consequently \(b_2\) cannot lie in
the completed field-norm image in A.159.2.
\end{proof}

\paragraph{Reproducible finite verification.}
The local PARI run is only a source of candidate coordinates.
\path{code/two_inert_relative_unit_repair_certificate.py}
checks all unit and relative-norm identities by exact rational arithmetic;
\path{code/two_inert_unit_three_saturation_certificate.py}
checks every entry of A.159.4 and all eighty-one exponent combinations.
The complete rational coordinates, integral characteristic polynomials,
residue values and program outputs are retained alongside those files.
The separate $S$-unit divisor certificate checks A.159.9--11 and the
norm-zero lattice.  No unverified class-group presentation is used.
